% Chapter 4 Introduction (untitled opening paragraphs)
% Funnel structure: broad context → GCE anomaly → debate → chapter roadmap
\aure{Review this introduction when I finalize the content of the chapters...}
\red{The preceding chapters established the theoretical basis for indirect dark matter detection through gamma rays (Chapter~1), described the instrumental and astrophysical landscape of the Fermi Large Area Telescope (Chapter~2), and introduced the statistical and machine-learning methods used to disentangle overlapping emission components (Chapter~3). We now turn to the most persistent anomaly in GeV gamma-ray observations: the Galactic Center Excess (GCE).}

Since 2009, analyses of Fermi-LAT data have consistently revealed a bright, statistically significant excess of gamma-ray emission from the region surrounding the Galactic Center~\cite{Goodenough:2009gk, Hooper:2010mq, Daylan:2014rsa, Calore:2014xka}. This signal peaks at energies between roughly 1 and 3~GeV, extends out to at least 10 degrees from Sgr~A$^*$, and persists across all existing models of the astrophysical diffuse emission~\cite{Calore:2014xka}. Its spectrum, angular distribution, and overall intensity are each consistent with the signal predicted from the annihilation of $\sim 40$--$70$~GeV dark matter particles distributed according to a contracted Navarro--Frenk--White profile~\cite{Daylan:2014rsa, Cirelli:2024ssz}. The best-fit annihilation cross section, $\langle \sigma v \rangle \sim (1$--$2) \times 10^{-26}$~cm$^3$/s, is compatible with the canonical thermal relic value~\cite{Daylan:2014rsa}.

The leading astrophysical alternative attributes the excess to a large, unresolved population of millisecond pulsars (MSPs) in the Inner Galaxy~\cite{Abazajian:2010zy}. MSPs produce gamma-ray emission with a spectral shape similar to the GCE, and early statistical analyses appeared to support a point-source origin~\cite{Lee:2015fea, Bartels:2015aea}. However, the robustness of these results has since been questioned~\cite{Leane:2019xiy, Leane:2020pfc}, and the MSP hypothesis faces independent observational challenges: few pulsars have been detected in the Inner Galaxy, and their expected number based on known luminosity functions substantially exceeds the handful of candidates identified in current catalogs~\cite{Hooper:2016rap, Holst:2024fvb, Amerio:2024qor}.

After more than fifteen years, the origin of the GCE remains an open question. This chapter traces the observational history of the excess (Section~\ref{sec:4.2}), examines the competing dark matter and millisecond pulsar interpretations together with the counter-arguments against each (Section~\ref{sec:4.3}), and describes the systematics-driven impasse that has stalled progress (Section~\ref{sec:4.4}). We then motivate an independent, orthogonal approach: constraining the MSP luminosity function through observations of the Milky Way's globular cluster system (Section~\ref{sec:4.5}). Section~\ref{sec:4.6} then presents this analysis in full, reproducing the published study of MSP gamma-ray luminosity functions in globular clusters and its implications for the GCE.
